% !TEX program = xelatex
\documentclass[journal]{IEEEtran}
\usepackage{amsmath,amssymb}
\usepackage{graphicx}
\usepackage{booktabs}
\usepackage{array}
\usepackage{cite}
\begin{document}

\title{How Long Should the Measurement Observe? Observation Span as a First-Order Design Variable in Aero-Engine Remaining-Life Estimation}

\author{Anonymous}

\maketitle

\begin{abstract}
Remaining useful life (RUL) predictors for aero-engines are compared by pooled regression error, a convention that fixes the observation window at 50 cycles without testing that choice. We treat the span as a measurement design variable and report what it is worth. On the windows where a longer history exists, a 100-cycle window lowers the pooled error from 11.03 to 7.50 cycles, and a 150-cycle window reaches 4.97 where it applies. Because a long window is only available once the engine has been observed long enough, we schedule the span to the available history and add two devices: an engine-level mixup that acts as a variance-reducing regulariser, and a pre-registered guard that falls back to the short-window model when the two models disagree by more than 50 cycles.

Across two seeds and five engine-disjoint folds, the complete recipe reduces the pooled RMSE by 10.00 percent on FD002 and 9.44 percent on FD004 relative to the prior-augmented baseline, with a stratified interval of 9.72 percent and bounds [1.55, 1.04] cycles. All four metrics improve on both fleets. On the two single-condition fleets the same recipe is 0.79 and 2.40 percent better. Three controls attribute the span gain to the span rather than to specialisation, parameter count or model count. We also report nine model-side interventions, four of them with zero additional parameters, none of which helps the prior arm: on this benchmark the return is in how the measurement is taken, not in what is added to the model.

\end{abstract}

\begin{IEEEkeywords}
Remaining useful life, measurement design, observation window, C-MAPSS, regularisation, model selection.
\end{IEEEkeywords}

\section{Introduction}

Remaining useful life is the quantity that maintenance decisions consume, and it is almost always reported as a single regression score on a benchmark. The score is computed over windows: a fixed-length slice of sensor history ending at the prediction time. The C-MAPSS literature has converged on 50 cycles \cite{saxena2008}, and the choice is rarely justified, because shortening or lengthening the window changes how many windows the evaluation contains and therefore the score itself.

This paper asks a measurement question rather than a modelling one: how long should the observation be, and what does the answer cost or buy? The question is well posed because a window is a measurement instrument. A 50-cycle window observes 50 samples of a degradation process whose time constant is hundreds of cycles; a 100-cycle window observes twice as much of the same process at the price of being unavailable for the first 100 cycles of the engine's life.

We answer the question on FD002 and FD004, the two six-condition fleets, under a protocol frozen before the experiments: five engine-disjoint outer folds, a full training budget with no epoch selection, a 42.3k-parameter backbone, and four metrics with engine-clustered intervals. Section~\ref{sec:method} states the method, Section~\ref{sec:protocol} the protocol, Section~\ref{sec:results} the results, and Section~\ref{sec:disc} the limits.

Our contributions are three. First, the observation span is a first-order design variable: scheduling it to the available history is worth more than any model-side change we were able to make. Second, the lifetime prior that individualisation methods inject is worth one to three percent, and an engine-level mixup with no additional parameters recovers the same amount, which suggests the prior acts as a regulariser rather than as information. Third, nine model-side interventions, four of them parameter-free, fail to improve the prior arm, which closes the explanation that the right architectural change had simply not been found.

\section{Related Work}

\subsection{Benchmarks and the window convention}
C-MAPSS \cite{saxena2008} remains the standard aero-engine RUL benchmark. Its fleets differ in the number of operating conditions and fault modes: FD001 and FD003 are single-condition, FD002 and FD004 combine six conditions. Nearly all reported results use a sliding window whose length is fixed and rarely ablated; the choice is inherited rather than tested. Our experiments make the span the independent variable.

\subsection{Individualisation and lifetime priors}
Qu \emph{et al.} \cite{qu2026} build a condition-aware spatio-temporal network in which operating settings and sensor measurements interact through a feature-interaction module and a hierarchical graph attention network. Kapoor \emph{et al.} \cite{kapoor2026} introduce a phase-aware normalisation with delta features for multi-regime engines. Both inject information about the individual engine beyond the raw window. We keep the backbone fixed and ask what that information is worth, and whether it can be obtained without a lifetime estimate.

\subsection{Uncertainty and preprocessing audits}
Piao \emph{et al.} \cite{piao2025} convert a point predictor into a conformal interval with coverage guarantees. Our intervals are engine-clustered bootstrap percentiles and carry no coverage guarantee; we do not describe our RMSE as an instrument uncertainty. Gupta \cite{gupta2026} audits operating-condition normalisation in C-MAPSS and shows that per-condition normalisation lowers error in every matched multi-regime cell; we take that as an established premise and freeze the fit scope in our protocol rather than presenting it as a contribution.

\subsection{Regularisation}
Mixup \cite{zhang2018} interpolates pairs of training examples and their labels, and is understood as a data-dependent regulariser that smooths the decision function. We apply it at the engine level, which for a fleet of a few hundred engines is the level at which the effective sample size is set.

\section{Method}\label{sec:method}

\subsection{Backbone}
Each observation is a window over 14 sensors plus three operating settings. The backbone is a channel-independent transformer with 42\,306 parameters ($d_{\mathrm{model}}=64$, one layer, dropout 0.15) that receives a positional encoding of the cycle coordinate, a condition-normalised sensor block, and a per-engine early-life reference state. Training uses AdamW (learning rate $10^{-3}$, weight decay $10^{-4}$), cosine annealing over 80 epochs, batch size 256, and the loss
\begin{equation}
\mathcal{L}=\mathrm{MSE}+\tfrac{1}{2}\,\mathrm{asym}(e),\qquad
\mathrm{asym}(e)=\begin{cases}1.8\,e^2, & e>0\\ e^2, & e\le 0,\end{cases}
\end{equation}
with $e=\hat{y}-y$ and the RUL target capped at 125 cycles.

\subsection{Observation-span schedule}
Two models are trained with the same budget and the same recipe, one on 50-cycle windows and one on 100-cycle windows. At inference the model is chosen by the available history: a window ending before 100 observed cycles uses the 50-cycle model, a window ending at or after 100 cycles uses the 100-cycle model. The rule is deterministic and uses only information available at the prediction time.

\subsection{Engine-level mixup}
Two windows from different engines are combined as $\lambda x_i+(1-\lambda)x_j$ with $\lambda\sim\mathrm{Beta}(0.2,0.2)$, and the label, the auxiliary inputs and the lifetime prior are combined the same way. The device adds no parameters. Mixing every channel consistently matters: mixing the window and the label but keeping the original prior turns a gain into a loss (Section~\ref{sec:negative}).

\subsection{Tail guard}
The two models can disagree. When they disagree by more than 50 cycles we keep the short-window prediction. The value is about 3.7 times the typical FD004 error; we report 30, 50 and 75 cycles as a sensitivity check. The rule was written into the protocol before any guarded result was computed.

\section{Protocol}\label{sec:protocol}

Folds are engine-disjoint: no engine contributes windows to both training and validation. The training budget is full for every arm (80 epochs over all available windows, final epoch used, no epoch selection), which matters because it removes early stopping as a confound and makes additional capacity strictly more expensive.

Four metrics are reported: RMSE, MAE, the asymmetric score $\exp(e/10)-1$ for $e>0$ and $\exp(-e/13)-1$ otherwise, and bias. Engines are weighted equally: per-engine metrics are computed first and then averaged, so a long-lived engine does not count more than a short-lived one. Intervals are engine-clustered paired bootstrap with 5000 resamples; the resampling unit is the engine, and all seeds of a resampled engine move together.

Two baselines are reported throughout. The plain backbone (PE) carries no lifetime prior. The prior arm injects $(\mu_T,\sigma_T)$ through a zero-initialised linear map, where $\mu_T$ is a HistGradientBoosting regressor over 88 features fitted on training engines only with an inner four-fold split by engine.

\section{Results}\label{sec:results}

\input{paper_tables/main_table.tex}

Table~\ref{tab:main} reports the four metrics. On FD002 the complete recipe reaches 11.871 against 13.189 for the prior arm, a reduction of 10.00 percent; on FD004 it reaches 12.173 against 13.440, a reduction of 9.44 percent. The stratified interval is 9.72 percent with bounds [1.55, 1.04] cycles, and the relative interval is [11.74, 7.78] percent. All four metrics improve on both fleets.

\subsection{The span is worth more than the window it replaces}
On FD002 fold 2, restricting attention to windows whose history is long enough to support both models (at least 100 observed cycles), a 100-cycle window lowers the pooled error from 11.03 to 7.50 cycles; on windows with at least 150 observed cycles, 50, 100 and 150-cycle windows give 7.04, 5.09 and 4.97. The comparison holds the window set fixed, so it is not an artefact of a changing evaluation set.

Three controls separate the span from competing explanations. A late-window specialist trained at the same 50-cycle window on the same late windows is worth 3.5 percent, so specialisation alone is small. A span control with identical parameter count, which samples 50 positions every second cycle over 100 cycles, is worth 12.8 percent, so both the span and the temporal resolution contribute. A three-seed ensemble of the 50-cycle reference, which has the same model count as the schedule, leaves the schedule ahead by 5.8 percent, so the gain is not a model-count effect.

\subsection{The guard removes a tail that the score amplifies}
The schedule alone is not stable on the asymmetric score. In one seed and one fold a single engine is over-estimated by about 96 cycles; because the score grows exponentially on the over-estimation side, that engine alone contributes 14695 to a fleet mean of 313, against 23 for the reference. The maximum absolute error on that fold is the same for both arms, 124.6 against 124.5; only the sign differs.

With the guard, FD004 moves from 2.92 to 3.89 percent and its asymmetric score from 37.76 to 8.68, below the reference value of 9.94. The gain is positive for every value in the sensitivity set {30, 50, 75}. The guard also repairs the plain backbone: FD004 moves from 2.03 percent with an interval crossing zero to 4.39 percent with bounds [0.85, 0.35].

\subsection{Replacing the prior with mixup}
The prior arm is not the strongest baseline on every fleet. On FD004 the plain backbone reaches an asymmetric score of 8.98 against 11.25 for the prior arm, so injecting the lifetime statistics helps the pooled error and hurts the tail. This suggests replacing the prior with a regulariser that acts on the model rather than on its input.

Engine-level mixup does that with no additional parameters. Alone, it brings the plain backbone to parity with the prior arm on FD002 (+0.3 percent for one seed, 3.1 percent for the other) and ahead of it on FD004 (1.7 and 6.0 percent), with the FD004 asymmetric score falling from 8.63 to 6.65 for the second seed. Averaged over the two seeds the gain is 1.4 percent on FD002 and 3.8 percent on FD004. The gain is larger for the seed whose initialisation is worse, which is what one expects from a regulariser that reduces variance.

Combined with the span schedule and the guard, the same model reaches the numbers in Table~\ref{tab:main}. The three devices act on different parts of the problem and the gains do not simply add: mixup alone is worth 1.4 percent on FD002, the span schedule with the guard 7.2 percent, and together 10.0 percent.

\subsection{Scope}
The complete recipe is better than the prior arm on all four fleets, but the size of the gain tracks the number of operating conditions. On the two single-condition fleets it is 0.79 and 2.40 percent better; on the two six-condition fleets it is 10.00 and 9.44 percent better.

Mixup alone shows the opposite pattern: on the single-condition fleets it is 1.04 and 4.67 percent worse than the prior arm, although on FD003 it remains 2.13 percent better than the plain backbone. The prior therefore carries information that nothing else supplies when the operating condition is fixed; the span schedule recovers that loss, and under six conditions the mixed-condition regularisation is worth more than the prior in the first place. We read the pattern as a scope statement rather than a caveat: the schedule exploits the longer history where that history is more informative, which is under varying conditions.

\subsection{What does not work}\label{sec:negative}
The span schedule changes how the measurement is taken. We tested nine model-side interventions with the same recipe, folds and metrics. Table~\ref{tab:neg} reports the single-fold screen against the 42k baseline on the same fold.

\begin{table}[t]
\centering
\caption{Model-side interventions, single-fold screen. Only one cell is clearly positive, and it is on the arm that does not carry the prior.}
\label{tab:neg}
\small
\begin{tabular}{lrrrr}
\toprule
Intervention & FD002 M0 & FD002 FULL & FD004 M0 & FD004 FULL \\
\midrule
Modules, +74k to +677k params & +7.4\% & +24.9\% & -- & -- \\
A: monotonicity loss & +4.3\% & not run & -- & -- \\
B: pinball loss & $-$3.9\% & +0.1\% & +5.3\% & +5.8\% \\
C: lifetime reparameterisation & +13.8\% & +9.9\% & -- & -- \\
D: engine-level mixup & $-$9.1\% & +0.8\% & -- & -- \\
DX: mixup on inputs only & +3.5\% & +4.5\% & -- & -- \\
E: cross-sensor attention & +6.4\% & -- & -- & -- \\
\bottomrule
\end{tabular}
\end{table}

Three structural conclusions follow. First, the prior arm is saturated: every intervention that helps the no-prior arm returns to zero or reverses on the prior arm, because the zero-initialised injection is itself a strong regulariser. Second, reducing degrees of freedom works on the no-prior arm while adding expressiveness never works, which is consistent with an effective sample size of about two hundred engines rather than thirty thousand windows. Third, mixup only works when every input channel is mixed consistently.

\section{Discussion}\label{sec:disc}

The study is simulation-based. The reported RMSE values are estimator performance metrics, not instrument uncertainties, and the intervals are descriptive engine-clustered bootstrap intervals rather than coverage-guaranteed statements. The official test engines have been used by this project before, so those numbers are post-hoc confirmations rather than a blind validation.

FD001 and FD003 are reported at a single seed; FD002 and FD004 at two. The training script initialises the network before resetting the random seed, so a single run is not bit-reproducible; multi-seed means are stable, and the conclusions rest on them rather than on any individual run.

The screening protocol deserves a caveat of its own. Our single-fold screens use roughly fifty engines, and the run-to-run noise of the frozen recipe is about three percent, which is the same order as the effects being screened. One arm looked like a 3.5 percent gain on two folds and settled at 0.3 percent on five; another moved from 15.7 percent on one fold to 2.4 percent on five. We therefore report the screens as direction-finding only, and the full matrices carry the conclusions.

Two further observations bear on how the results should be read. The plain backbone is a strong baseline, only 3.3 percent behind the prior-augmented model on FD002 and 1.1 percent behind on FD004, so the prior is worth one to three percent. And an engine-level mixup with no parameters reaches parity with it, which suggests the prior's contribution is regularisation rather than information.

\section{Conclusion}

The observation span is a first-order design variable in RUL measurement, and on C-MAPSS it is worth more than the model-side changes we were able to make. Scheduling the span to the available history, regularising with engine-level mixup, and guarding the tail gives 10.00 and 9.44 percent over the prior-augmented baseline on the two multi-condition fleets, with all four metrics improving and an engine-clustered stratified interval of [11.74, 7.78] percent. Nine model-side interventions, four of them parameter-free, fail to improve the prior arm, which closes the explanation that the right architectural change had not been found.

We recommend that RUL studies report the observation span as an explicit design choice, alongside the evaluation endpoint, the engine sample and the observability of any individualisation signal they inject.

\begin{thebibliography}{99}

\bibitem{saxena2008} A.~Saxena, K.~Goebel, D.~Simon, and N.~Eklund, ``Damage propagation modeling for aircraft engine run-to-failure simulation,'' in \emph{Proc. Int. Conf. Prognostics and Health Management}, 2008.

\bibitem{qu2026} Y.~Qu, X.~Yang, Y.~Luo, H.~Liu, Z.~Zhou, and L.~Guo, ``A condition-aware spatio-temporal network for robust aero-engine remaining useful life prediction,'' \emph{IEEE Trans. Instrum. Meas.}, 2026, doi: 10.1109/TIM.2026.3742092.

\bibitem{kapoor2026} R.~Kapoor, S.~Gupta, A.~Agarwal, K.~Uppal, S.~Juneja, and P.~Sharma, ``STARNet: Stacked transfer-aware for robust remaining useful life prediction for C-MAPSS multi-regime engines,'' \emph{IEEE Access}, 2026, doi: 10.1109/ACCESS.2026.3663754.

\bibitem{piao2025} Y.~Piao, C.~Huang, and F.~Tsung, ``CRULP: Reliable RUL estimation inspired by conformal prediction,'' \emph{IEEE Trans. Instrum. Meas.}, vol.~74, 2025, doi: 10.1109/TIM.2024.3522678.

\bibitem{gupta2026} A.~Gupta, ``Auditing operating-condition normalization in C-MAPSS remaining useful life evaluation,'' engrXiv preprint, 2026, doi: 10.31224/8145.

\bibitem{zhang2018} H.~Zhang, M.~Cisse, Y.~N.~Dauphin, and D.~Lopez-Paz, ``mixup: Beyond empirical risk minimization,'' in \emph{Proc. Int. Conf. Learning Representations}, 2018.

\end{thebibliography}

\end{document}